\documentclass{zhvt-classic}

% 使用本地源流明体（GenRyuMin2 TW）按文件路径加载，不依赖系统字体名
% 版框（内框／界栏／书口侧栏／鱼尾）改用与版心印章同一印泥棕色：
%   印章图片 yshs.png 实色 = RGB(135,68,52) = HTML 874434
\usepackage{xcolor}
\usepackage{tikz}
% sealbrown 双注册：xcolor 供 TikZ/pgf，L3 供 cls 版框
\definecolor{sealbrown}{HTML}{874434}
\ExplSyntaxOn
\color_set:nnn { sealbrown } { HTML } { 874434 }
\ExplSyntaxOff

\zhvtset{
  grid color       = sealbrown,
  main font        = GenRyuMin2TW-SB,
  main font path   = C:/Users/Administrator/texmf/fonts/opentype/public/GenRyuMin2TW/,
jiazhu font      = GenRyuMin2TW-M,
  jiazhu font path = C:/Users/Administrator/texmf/fonts/opentype/public/GenRyuMin2TW/,
  % 页码框容纳 4 字：正常页「九十七」+前置空 1 格=4 行恰好；
  % 100+ 页「一百零一」共 5 行超出，末字「一/二/三」被框裁掉（2026-10-04 实测），故放宽到 5
  shukou page number chars = 5,
  % 书口书名框高从默认 6 字减到 4 字，把省出的 2 字高拨给 chapter 框，
  % 使 6 字书名+卷名放进 9 字高的 chapter 框后不再压线（2026-10-05）
  shukou title chars = 4,
}

\title{史记}[齐太公世家]
\author{汉·司马迁}
\maker{闲叟漫记·重製}


% ===== 索隐述赞样式（2026-11-08 四次定稿 · 愚大人亲手改定参数）=====
% ⚠ 竖排前提：自定义字体家族必须带 RawFeature={vertical:+vert}；
%   simhei.ttf 加 +vhal 会整体加载失败静默回退，只写 +vert；simfang 同理。
\newCJKfontfamily{\shuzanfont}[Path=C:/Windows/Fonts/,RawFeature={vertical:+vert}]{simhei.ttf}
\newCJKfontfamily{\shuzanbody}[Path=C:/Windows/Fonts/,RawFeature={vertical:+vert}]{simfang.ttf}
% 题头：152x30pt 圆角 6pt 棕底白字；baseline (0,2) 使整框落入界行格子不压版框线；
% 白字节点 x 不补偿、y +8.2pt（像素级墨迹实测 dx/dy≈0，L/R 留白约 20pt）。
\newcommand{\shuzanhead}[1]{%
  \begin{tikzpicture}[baseline={(0,2)}]
    \fill[sealbrown, rounded corners=6pt]
      (-76pt,-15pt) rectangle (76pt,15pt);
    \node[text=white, inner sep=0pt, anchor=center] at (0pt,8.2pt)
      {\shuzanfont\fontsize{23}{23}\selectfont #1};
  \end{tikzpicture}%
}
% ★ 修正书口卷名位置：chapter 框起始位置改回 cls 默认值，避免全局影响。
%   只在每个卷首页通过 \juan 命令里的空命令微调该页位置（安全，不会跨页污染）。
\ExplSyntaxOn
\AtBeginDocument{
  \dim_set:Nn \l__zhvt_shukou_chapter_title_start_dim
    { \l__zhvt_shukou_title_height_dim + \l__zhvt_font_size_dim * 2 }
  \dim_set:Nn \l__zhvt_shukou_page_number_start_dim
    { \l__zhvt_shukou_chapter_title_start_dim + \l__zhvt_shukou_chapter_title_height_dim }
}
\ExplSyntaxOff

\begin{document}

% ★★ 版心（书口）排布覆写：把书名放到鱼尾之下（2026-10-04）
%
%   cls 默认版心自上而下：书名（title 静态框，在鱼尾之上）→ 鱼尾 → 篇名（chapter 动态框）→ 页码
%   本模板改为：          鱼尾（位置不动）→ 书名 → 篇名 → 页码分割线 → 页码
%
%   手法：把 cls 放在 title 静态框里的书名，改放进 chapter 动态框（鱼尾之下）；
%         title 静态框留空。书名取 \l__zhvt_title_str，即 \title{...} 的第一参数，
%         所以换书时只改 \title，不必在此重复写书名。
%
%   ⚠ 三条必须遵守的写法（2026-10-04 实测）：
%     ① 汉字必须用 \hbox{...} 包住。flowfram 会把框内容装进一个宽 14.4mm（= 一个 line_sep
%        = 书口栏宽）的纵向 minipage 并按该宽度折行，而整页靠 PDF /Rotate 90 显示、
%        框内文字本身不变向 —— 不包 \hbox 时「围炉夜话」4 字会被逐字折成 4 行，
%        在版面上就成了横排 4 字并横向溢出版框。
%     ② 竖排下**一个 \hbox 就是一整列**。书名与篇名要同处一列，必须放进同一个 \hbox；
%        分成两个 \hbox 时，后一个会被当成新的一列、甩到隔壁栏去
%        （实测「序」跑到书口栏外侧 14.4mm 的正文列上）。中间的 \quad 在框内是横向空白，
%        显示出来正好是书名与篇名之间的纵向间隔。
%     ③ ★ 内容盒要写成 **与框等宽的定宽盒** `\hbox to \l__zhvt_shukou_chapter_title_height_dim{…\hfil}`。
%        实测：chapter 框是把内容盒沿纵向**居中**摆放的，内容宽度一变，书名就跟着漂：
%        自然宽度的 \hbox（书名 38.4mm）落在 x=157.2mm，加篇名后（57.6mm）落在 x=152.4mm，
%        目录页（无篇名）与正文页的书名就差半格。做成等宽盒后，无论有没有篇名，
%        内容固定从框顶起排，实测两页均落在 x=142.8mm，不再漂。
\ExplSyntaxOn
% 把书口 chapter 框（书名+篇名）起始位置下移 1 字高，避免 4 字卷名顶到鱼尾
% ⚠ cls 的 \__zhvt_calculate_dimensions: 在 \AtBeginDocument 里执行，会重算该 dim；
%   必须把覆写挂到 begin document 之后，否则会被二次覆盖（2026-10-05 实测）。
% (旧方案已移到导言区，直接修改 chapter 框起始位置)
% 卷标题命令 #1=书口篇名(鱼尾) #2=正文卷名（含卷序，卷名内空格用 ~）
%
% ★★ 目录页两项修订（2026-10-05，愚大人指出）★★
%   ① 页码起点统一对齐：
%      原状 —— 目录项写成 \par\hspace{4em}卷名\hspace{4em}页码，页码起点 = 卷名实宽 + 4em。
%      卷名实宽随字数变：「卷一~梁惠王上」6 字 +~、「卷七~离娄上」5 字 +~，
%      差 1 个字位 → 卷七~卷十 4 条的页码整体比其余 10 条高 27.21pt（一个字），看着上下不齐。
%      改法 —— 卷名补全角空格（宽 = 1 字 = 27.215pt）补足到 **7 字符位**再写进目录，
%      14 条页码起点遂全部落在同一条竖线上（起点对齐；末字随页码字数自然参差）。
%      ⚠ 只对目录里的副本 \l_tmpc_tl 补，正文卷名仍用 \l_tmpa_tl 原样 —— 正文像素不受影响。
%   ② 卷名整体上移半字位：前导 \hspace 由 4em 改 3.5em。
%      实测原状目录块在版框内偏下（天头留白 139.6pt / 地脚留白 110.7pt），
%      上移 0.5em = 13.6pt 后上下基本均分，视觉居中。
\NewDocumentCommand { \juan } { mm } {
	\clearpage
	\tl_set:Nn \l__zhvt_chapter_title_tl { #1 }
	\tl_set:Nn \l_tmpa_tl { #2 }
	\tl_set_eq:NN \l_tmpc_tl \l_tmpa_tl
	\int_step_inline:nn { 7 - \str_count:n { #2 } }
	  { \tl_put_right:Nn \l_tmpc_tl { 　 } }
	\tl_set:Nx \l_tmpb_tl { \zhnum { page } }
	\tl_put_left:Nn \l_tmpb_tl { \hspace { 4em } }
	\tl_put_left:Nn \l_tmpb_tl { \l_tmpc_tl }
	\tl_put_left:Nn \l_tmpb_tl { \par \hspace { 3.5em } }
	\addtocontents { toc } { \l_tmpb_tl }
	\skip_horizontal:n { \l__zhvt_font_size_dim * 4 }
	\l_tmpa_tl
	\par
	\thispagestyle { empty }
}
\cs_set_protected:Npn \__zhvt_set_shukou_content: {
	% ★ 鱼尾上方留白（title 静态框）嵌入闲叟钤印：
	%   整页 /Rotate 90，故此处把图用 \rotatebox{90} 转正为竖章；
	%   尺寸卡在版心条幅内（书内约 7mm 宽 × 33.8mm 高，留余量不压鱼尾）。
	\setstaticcontents * { title } {
		\vbox_to_ht:nn { \l__zhvt_baseline_skip_dim } {
			\vskip 14.2pt
			\hbox_to_wd:nn { \l__zhvt_shukou_title_height_dim } {
				\hfil
				\rotatebox [ origin = c ] { 90 } { \includegraphics [ width = 7mm , height = 33.8mm ] { yshs.png } }
				\hfil
			}
		}
	}
% ★ 版心（书口）统一改法（2026-11-16 愚大人定）：18pt＋框内居中＋书口文字同版框色
%   ① 原 27.2pt 时 7 字长篇名（如「廉颇蔺相如列传」）末字溢出 chapter 框、压住页码；
%      改 \fontsize{18pt}{18pt}，前后 \hfil 令字串在框内居中（不随篇名字数漂移）。
%   ② \vbox{\vskip 8.2pt} 定量下移，抵消缩字号后栏宽方向 8.2pt 上偏（像素实测 y 中心 543.3 = 原版）。
%      ⚠ \raisebox / \vspace* 在 flowfram 框内无效；换字号须同步换补偿：16pt→10.1、19pt→7.4。
%   ③ chapter 与 page~number 均置 \color{sealbrown}（HTML 874434），与版框／界行同色；正文仍黑。
%   ⚠ 页码字位、钤印、鱼尾、界行、正文一律未动；五部合订本同此口径。
	\setdynamiccontents * { chapter }
	{ \color{sealbrown} \vbox { \vskip 8.2pt \hbox to \l__zhvt_shukou_chapter_title_height_dim
		{ \hfil \fontsize{18pt}{18pt}\selectfont \l__zhvt_title_str \quad \l__zhvt_chapter_title_tl \hfil } } }
	\setdynamiccontents * { page~number } { \color{sealbrown} \quad \zhnum { page } }
}
\__zhvt_set_shukou_content:
\ExplSyntaxOff

\maketitle

% 单篇，无多卷目录，不渲染目录页

\juan{齐太公世家}{卷三十二~齐太公世家第二}
\hspace*{2em}太公望卜尚者。东海上人。其先功尝为四岳。佐禹平水土甚有功。虞夏之际封于卜。或封于申。庶姜氏。夏商之时。申。卜或封枝庶子孙。或为庶人。尚其后苗裔也。本庶姜氏。从其封庶。故曰卜尚。\\
\hspace*{2em}卜尚盖尝穷困。年老矣。以渔钓见周西伯。西伯将出猎。卜之。曰﹁所获非龙非彨。非虎非罴。所获霸王之辅﹂。于是周西伯猎。果遇太公于渭之阳。与语大说。曰。﹁自吾先君太公曰﹃当有圣人适周。周以兴﹄。子真是邪。吾太公望子久矣。﹂故号之曰﹁太公望﹂。载与俱归。立为师。\\
\hspace*{2em}或曰。太公博闻。尝事纣。纣无道。去之。游说诸侯。无所遇。而卒西归周西伯。或曰。卜尚处士。隐海滨。周西伯拘羑里。散宜生。闳夭素知而招卜尚。卜尚亦曰﹁吾闻西伯贤。又善养老。盍往焉﹂。三人者为西伯求美女奇物。献之于纣。以赎西伯。西伯得以出。反国。言卜尚所以事周虽异。然要之为文武师。\\
\hspace*{2em}周西伯昌之脱羑里归。与卜尚阴谋修德以倾商政。其事多兵权与奇计。故后世之言兵及周之阴权皆宗太公为本谋。周西伯政平。及断虞芮之讼。而诗人称西伯受命曰文王。伐崇。密须。犬夷。大作丰邑。天下三分。其二归周者。太公之谋计居多。\\
\hspace*{2em}文王崩。武王即位。九年。欲修文王业。东伐以观诸侯集否。师行。师尚父左杖黄钺。右把白旄以誓。曰。﹁苍兕苍兕。总尔众庶。与尔舟楫。后至者斩。﹂遂至盟津。诸侯不期而会者八百诸侯。诸侯皆曰。﹁纣可伐也。﹂武王曰。﹁未可。﹂还师。与太公作此太誓。\\
\hspace*{2em}居二年。纣杀王子比干。囚箕子。武王将伐纣。卜龟兆。不吉。风雨暴至。群公尽惧。唯太公强之劝武王。武王于是遂行。十一年正月甲子。誓于牧野。伐商纣。纣师败绩。纣反走。登鹿台。遂追斩纣。明日。武王立于社。群公奉明水。卫康叔封布衣席。师尚父牵牲。史佚策祝。以告神讨纣之罪。散鹿台之钱。发钜桥之粟。以振贫民。封比干墓。释箕子囚。迁九鼎。修周政。与天下更始。师尚父谋居多。\\
\hspace*{2em}于是武王已平商而王天下。封师尚父于齐营丘。东就国。道宿行迟。逆旅之人曰。﹁吾闻时难得而易失。客寝甚安。殆非就国者也。﹂太公闻之。夜衣而行。犁明至国。莱侯来伐。与之争营丘。营丘边莱。莱人。夷也。会纣之乱而周初定。未能集远方。是以与太公争国。\\
\hspace*{2em}太公至国。修政。因其俗。简其礼。通商工之业。便鱼盐之利。而人民多归齐。齐为大国。及周成王少时。管蔡作乱。淮夷畔周。乃使召康公命太公曰。﹁东至海。西至河。南至穆陵。北至无棣。五侯九伯。实得征之。﹂齐由此得征伐。为大国。都营丘。\\
\hspace*{2em}盖太公之卒百有余年。子丁公卜伋立。丁公卒。子乙公得立。乙公卒。子癸公慈母立。癸公卒。子哀公不辰立。\\
\hspace*{2em}哀公时。纪侯谮之周。周烹哀公而立其弟静。是为胡公。胡公徙都薄姑。而当周夷王之时。\\
\hspace*{2em}哀公之同母少弟山怨胡公。乃与其党率营丘人袭攻杀胡公而自立。是为献公。献公元年。尽逐胡公子。因徙薄姑都。治临菑。\\
\hspace*{2em}九年。献公卒。子武公寿立。武公九年。周厉王出奔。居彘。十年。王室乱。大臣行政。号曰﹁共和﹂。二十四年。周宣王初立。\\
\hspace*{2em}二十六年。武公卒。子厉公无忌立。厉公暴虐。故胡公子复入齐。齐人欲立之。乃与攻杀厉公。胡公子亦战死。齐人乃立厉公子赤为君。是为文公。而诛杀厉公者七十人。\\
\hspace*{2em}文公十二年卒。子成公脱立。成公九年卒。子庄公购立。\\
\hspace*{2em}庄公二十四年。犬戎杀幽王。周东徙雒。秦始列为诸侯。五十六年。晋弑其君昭侯。\\
\hspace*{2em}六十四年。庄公卒。子厘公禄甫立。\\
\hspace*{2em}厘公九年。鲁隐公初立。十九年。鲁桓公弑其兄隐公而自立为君。\\
\hspace*{2em}二十五年。北戎伐齐。郑使太子忽来救齐。齐欲妻之。忽曰。﹁郑小齐大。非我敌。﹂遂辞之。\\
\hspace*{2em}三十二年。厘公同母弟夷仲年死。其子曰公孙无知。厘公爱之。令其秩服奉养比太子。\\
\hspace*{2em}三十三年。厘公卒。太子诸儿立。是为襄公。\\
\hspace*{2em}襄公元年。始为太子时。尝与无知斗。及立。绌无知秩服。无知怨。\\
\hspace*{2em}四年。鲁桓公与夫人如齐。齐襄公故尝私通鲁夫人。鲁夫人者。襄公女弟也。自厘公时嫁为鲁桓公妇。及桓公来而襄公复通焉。鲁桓公知之。怒夫人。夫人以告齐襄公。齐襄公与鲁君饮。醉之。使力士彭生抱上鲁君车。因拉杀鲁桓公。桓公下车则死矣。鲁人以为让。而齐襄公杀彭生以谢鲁。\\
\hspace*{2em}八年。伐纪。纪迁去其邑。\\
\hspace*{2em}十二年。初。襄公使连称。管至父戍葵丘。瓜时而往。及瓜而代。往戍一岁。卒瓜时而公弗为发代。或为请代。公弗许。故此二人怒。因公孙无知谋作乱。连称有从妹在公宫。无宠。使之间襄公。曰﹁事成以女为无知夫人﹂。冬十二月。襄公游姑棼。遂猎沛丘。见彘。从者曰﹁彭生﹂。公怒。射之。彘人立而啼。公惧。坠车伤足。失屦。反而鞭主屦者茀三百。茀出宫。而无知。连称。管至父等闻公伤。乃遂率其众袭宫。逢主屦茀。茀曰。﹁且无入惊宫。惊宫未易入也。﹂无知弗信。茀示之创。乃信之。待宫外。令茀先入。茀先入。即匿襄公户间。良久。无知等恐。遂入宫。茀反与宫中及公之幸臣攻无知等。不胜。皆死。无知入宫。求公不得。或见人足于户间。发视。乃襄公。遂弑之。而无知自立为齐君。\\
\hspace*{2em}桓公元年春。齐君无知游于雍林。雍林人尝有怨无知。及其往游。雍林人袭杀无知。告齐大夫曰。﹁无知弑襄公自立。臣谨行诛。唯大夫更立公子之当立者。唯命是听。﹂\\
\hspace*{2em}初。襄公之醉杀鲁桓公。通其夫人。杀诛数不当。淫于妇人。数欺大臣。群弟恐祸及。故次弟纠奔鲁。其母鲁女也。管仲。召忽傅之。次弟小白奔莒。鲍叔傅之。小白母。卫女也。有宠于厘公。小白自少好善大夫高傒。及雍林人杀无知。议立君。高。国先阴召小白于莒。鲁闻无知死。亦发兵送公子纠。而使管仲别将兵遮莒道。射中小白带钩。小白详死。管仲使人驰报鲁。鲁送纠者行益迟。六日至齐。则小白已入。高傒立之。是为桓公。\\
\hspace*{2em}桓公之中钩。详死以误管仲。已而载温车中驰行。亦有高。国内应。故得先入立。发兵距鲁。秋。与鲁战于干时。鲁兵败走。齐兵掩绝鲁归道。齐遗鲁书曰。﹁子纠兄弟。弗忍诛。请鲁自杀之。召忽。管仲雠也。请得而甘心醢之。不然。将围鲁。﹂鲁人患之。遂杀子纠于笙渎。召忽自杀。管仲请囚。桓公之立。发兵攻鲁。心欲杀管仲。鲍叔牙曰。﹁臣幸得从君。君竟以立。君之尊。臣无以增君。君将治齐。即高傒与叔牙足也。君且欲霸王。非管夷吾不可。夷吾所居国国重。不可失也。﹂于是桓公从之。乃详为召管仲欲甘心。实欲用之。管仲知之。故请往。鲍叔牙迎受管仲。及堂阜而脱桎梏。斋祓而见桓公。桓公厚礼以为大夫。任政。\\
\hspace*{2em}桓公既得管仲。与鲍叔。隰朋。高傒修齐国政。连五家之兵。设轻重鱼盐之利。以赡贫穷。禄贤能。齐人皆说。\\
\hspace*{2em}二年。伐灭郯。郯子奔莒。初。桓公亡时。过郯。郯无礼。故伐之。\\
\hspace*{2em}五年。伐鲁。鲁将师败。鲁庄公请献遂邑以平。桓公许。与鲁会柯而盟。鲁将盟。曹沫以匕首劫桓公于坛上。曰。﹁反鲁之侵地。﹂桓公许之。已而曹沫去匕首。北面就臣位。桓公后悔。欲无与鲁地而杀曹沫。管仲曰。﹁夫劫许之而倍信杀之。愈一小快耳。而弃信于诸侯。失天下之援。不可。﹂于是遂与曹沫三败所亡地于鲁。诸侯闻之。皆信齐而欲附焉。七年。诸侯会桓公于甄。而桓公于是始霸焉。\\
\hspace*{2em}十四年。陈厉公子完。号敬仲。来奔齐。齐桓公欲以为卿。让。于是以为工正。田成子常之功也。\\
\hspace*{2em}二十三年。山戎伐燕。燕告急于齐。齐桓公救燕。遂伐山戎。至于孤竹而还。燕庄公遂送桓公入齐境。桓公曰。﹁非天子。诸侯相送不出境。吾不可以无礼于燕。﹂于是分沟割燕君所至与燕。命燕君复修召公之政。纳贡于周。如成康之时。诸侯闻之。皆从齐。\\
\hspace*{2em}二十七年。鲁愍公母曰哀姜。桓公女弟也。哀姜淫于鲁公子庆父。庆父弑愍公。哀姜欲立庆父。鲁人更立厘公。桓公召哀姜。杀之。\\
\hspace*{2em}二十八年。卫文公有狄乱。告急于齐。齐率诸侯城楚丘而立卫君。\\
\hspace*{2em}二十九年。桓公与夫人蔡姬戏船中。蔡姬习水。荡公。公惧。止之。不止。出船。怒。归蔡姬。弗绝。蔡亦怒。嫁其女。桓公闻而怒。兴师往伐。\\
\hspace*{2em}三十年春。齐桓公率诸侯伐蔡。蔡溃。遂伐楚。楚成王兴师问曰。﹁何故涉吾地。﹂管仲对曰。﹁昔召康公命我先君太公曰。﹃五侯九伯。若实征之。以夹辅周室。﹄赐我先君履。东至海。西至河。南至穆陵。北至无棣。楚贡包茅不入。王祭不具。是以来责。昭王南征不复。是以来问。﹂楚王曰。﹁贡之不入。有之。寡人罪也。敢不共乎。昭王之出不复。君其问之水滨。﹂齐师进次于陉。夏。楚王使屈完将兵扞齐。齐师退次召陵。桓公矜屈完以其众。屈完曰。﹁君以道则可。若不。则楚方城以为城。江。汉以为沟。君安能进乎。﹂乃与屈完盟而去。过陈。陈袁涛涂诈齐。令出东方。觉。秋。齐伐陈。是岁。晋杀太子申生。\\
\hspace*{2em}三十五年夏。会诸侯于葵丘。周襄王使宰孔赐桓公文武胙。彤弓矢。大路。命无拜。桓公欲许之。管仲曰﹁不可﹂。乃下拜受赐。秋。复会诸侯于葵丘。益有骄色。周使宰孔会。诸侯颇有叛者。晋侯病。后。遇宰孔。宰孔曰。﹁齐侯骄矣。弟无行。﹂从之。是岁。晋献公卒。里克杀傒齐。卓子。秦穆公以夫人入公子夷吾为晋君。桓公于是讨晋乱。至高梁。使隰朋立晋君。还。\\
\hspace*{2em}是时周室微。唯齐。楚。秦。晋为强。晋初与会。献公死。国内乱。秦穆公辟远。不与中国会盟。楚成王初收荆蛮有之。夷狄自置。唯独齐为中国会盟。而桓公能宣其德。故诸侯宾会。于是桓公称曰。﹁寡人南伐至召陵。望熊山。北伐山戎。离枝。孤竹。西伐大夏。涉流沙。束马悬车登太行。至卑耳山而还。诸侯莫违寡人。寡人兵车之会三。乘车之会六。九合诸侯。一匡天下。昔三代受命。有何以异于此乎。吾欲封泰山。禅梁父。﹂管仲固谏。不听。乃说桓公以远方珍怪物至乃得封。桓公乃止。\\
\hspace*{2em}三十八年。周襄王弟带与戎。翟合谋伐周。齐使管仲平戎于周。周欲以上卿礼管仲。管仲顿首曰。﹁臣陪臣。安敢。﹂三让。乃受下卿礼以见。三十九年。周襄王弟带来奔齐。齐使仲孙请王。为带谢。襄王怒。弗听。\\
\hspace*{2em}四十一年。秦穆公虏晋惠公。复归之。是岁。管仲。隰朋皆卒。管仲病。桓公问曰。﹁群臣谁可相者。﹂管仲曰。﹁知臣莫如君。﹂公曰。﹁易牙如何。﹂对曰。﹁杀子以适君。非人情。不可。﹂公曰。﹁开方如何。﹂对曰。﹁倍亲以适君。非人情。难近。﹂公曰。﹁竖刀如何。﹂对曰。﹁自宫以适君。非人情。难亲。﹂管仲死。而桓公不用管仲言。卒近用三子。三子专权。\\
\hspace*{2em}四十二年。戎伐周。周告急于齐。齐令诸侯各发卒戍周。是岁。晋公子重耳来。桓公妻之。\\
\hspace*{2em}四十三年。初。齐桓公之夫人三。曰王姬。徐姬。蔡姬。皆无子。桓公好内。多内宠。如夫人者六人。长卫姬。生无诡。少卫姬。生惠公元。郑姬。生孝公昭。葛嬴。生昭公潘。密姬。生懿公商人。宋华子。生公子雍。桓公与管仲属孝公于宋襄公。以为太子。雍巫有宠于卫共姬。因宦者竖刀以厚献于桓公。亦有宠。桓公许之立无诡。管仲卒。五公子皆求立。冬十月乙亥。齐桓公卒。易牙入。与竖刀因内宠杀群吏。而立公子无诡为君。太子昭奔宋。\\
\hspace*{2em}桓公病。五公子各树党争立。及桓公卒。遂相攻。以故宫中空。莫敢棺。桓公尸在床上六十七日。尸虫出于户。十二月乙亥。无诡立。乃棺赴。辛巳夜。敛殡。\\
\hspace*{2em}桓公十有余子。要其后立者五人。无诡立三月死。无谥。次孝公。次昭公。次懿公。次惠公。孝公元年三月。宋襄公率诸侯兵送齐太子昭而伐齐。齐人恐。杀其君无诡。齐人将立太子昭。四公子之徒攻太子。太子走宋。宋遂与齐人四公子战。五月。宋败齐四公子师而立太子昭。是为齐孝公。宋以桓公与管仲属之太子。故来征之。以乱故。八月乃葬齐桓公。\\
\hspace*{2em}六年春。齐伐宋。以其不同盟于齐也。夏。宋襄公卒。七年。晋文公立。\\
\hspace*{2em}十年。孝公卒。孝公弟潘因卫公子开方杀孝公子而立潘。是为昭公。昭公。桓公子也。其母曰葛嬴。\\
\hspace*{2em}昭公元年。晋文公败楚于城濮。而会诸侯践土。朝周。天子使晋称伯。六年。翟侵齐。晋文公卒。秦兵败于殽。十二年。秦穆公卒。\\
\hspace*{2em}十九年五月。昭公卒。子舍立为齐君。舍之母无宠于昭公。国人莫畏。昭公之弟商人以桓公死争立而不得。阴交贤士。附爱百庶。百庶说。及昭公卒。子舍立。孤弱。即与众十月即墓上弑齐君舍。而商人自立。是为懿公。懿公。桓公子也。其母曰密姬。\\
\hspace*{2em}懿公四年春。初。懿公为公子时。与丙戎之父猎。争获不胜。及即位。断丙戎父足。而使丙戎仆。庸职之妻好。公内之宫。使庸职骖乘。五月。懿公游于申池。二人浴。戏。职曰。﹁断足子。﹂戎曰。﹁夺妻者。﹂二人俱病此言。乃怨。谋与公游竹中。二人弑懿公车上。弃竹中而亡去。\\
\hspace*{2em}懿公之立。骄。民不附。齐人废其子而迎公子元于卫。立之。是为惠公。惠公。桓公子也。其母卫女。曰少卫姬。避齐乱。故在卫。\\
\hspace*{2em}惠公二年。长翟来。王子城父攻杀之。埋之于北门。晋赵穿弑其君灵公。\\
\hspace*{2em}十年。惠公卒。子顷公无野立。初。崔杼有宠于惠公。惠公卒。高。国畏其逼也。逐之。崔杼奔卫。\\
\hspace*{2em}顷公元年。楚庄王强。伐陈。二年。围郑。郑伯降。已复国郑伯。\\
\hspace*{2em}六年春。晋使郄克于齐。齐使夫人帷中而观之。郄克上。夫人笑之。郄克曰。﹁不是报。不复涉河。﹂归。请伐齐。晋侯弗许。齐使至晋。郄克执齐使者四人河内。杀之。八年。晋伐齐。齐以公子强质晋。晋兵去。十年春。齐伐鲁。卫。鲁。卫大夫如晋请师。皆因郄克。晋使郄克以车八百乘为中军将。士燮将上军。栾书将下军。以救鲁。卫。伐齐。六月壬申。与齐侯兵合靡笄下。癸酉。陈于鞍。逄丑父为齐顷公右。顷公曰。﹁驰之。破晋军会食。﹂射伤郄克。流血至履。克欲还入壁。其御曰。﹁我始入。再伤。不敢言疾。恐惧士卒。愿子忍之。﹂遂复战。战。齐急。丑父恐齐侯得。乃易处。顷公为右。车絓于木而止。晋小将韩厥伏齐侯车前。曰﹁寡君使臣救鲁。卫﹂。戏之。丑父使顷公下取饮。因得亡。脱去。入其军。晋郄克欲杀丑父。丑父曰。﹁代君死而见僇。后人臣无忠其君者矣。﹂克舍之。丑父遂得亡归齐。于是晋军追齐至马陵。齐侯请以宝器谢。不听。必得笑克者萧桐叔子。令齐东亩。对曰。﹁叔子。齐君母。齐君母亦犹晋君母。子安置之。且子以义伐而以暴为后。其可乎。﹂于是乃许。令反鲁。卫之侵地。\\
\hspace*{2em}十一年。晋初置六卿。赏鞍之功。齐顷公朝晋。欲尊王晋景公。晋景公不敢受。乃归。归而顷公弛苑囿。薄赋敛。振孤问疾。虚积聚以救民。民亦大说。厚礼诸侯。竟顷公卒。百庶附。诸侯不犯。\\
\hspace*{2em}十七年。顷公卒。子灵公环立。\\
\hspace*{2em}灵公九年。晋栾书弑其君厉公。十年。晋悼公伐齐。齐令公子光质晋。十九年。立子光为太子。高厚傅之。令会诸侯盟于钟离。二十七年。晋使中行献子伐齐。齐师败。灵公走入临菑。晏婴止灵公。灵公弗从。曰。﹁君亦无勇矣。﹂晋兵遂围临菑。临菑城守不敢出。晋焚郭中而去。\\
\hspace*{2em}二十八年。初。灵公取鲁女。生子光。以为太子。仲姬。戎姬。戎姬嬖。仲姬生子牙。属之戎姬。戎姬请以为太子。公许之。仲姬曰。﹁不可。光之立。列于诸侯矣。今无故废之。君必悔之。﹂公曰。﹁在我耳。﹂遂东太子光。使高厚傅牙为太子。灵公疾。崔杼迎故太子光而立之。是为庄公。庄公杀戎姬。五月壬辰。灵公卒。庄公即位。执太子牙于句窦之丘。杀之。八月。崔杼杀高厚。晋闻齐乱。伐齐。至高唐。\\
\hspace*{2em}庄公三年。晋大夫栾盈奔齐。庄公厚客待之。晏婴。田文子谏。公弗听。四年。齐庄公使栾盈间入晋曲沃为内应。以兵随之。上太行。入孟门。栾盈败。齐兵还。取朝歌。\\
\hspace*{2em}六年。初。棠公妻好。棠公死。崔杼取之。庄公通之。数如崔氏。以崔杼之冠赐人。待者曰。﹁不可。﹂崔杼怒。因其伐晋。欲与晋合谋袭齐而不得间。庄公尝笞宦者贾举。贾举复侍。为崔杼间公以报怨。五月。莒子朝齐。齐以甲戌飨之。崔杼称病不视事。乙亥。公问崔杼病。遂从崔杼妻。崔杼妻入室。与崔杼自闭户不出。公拥柱而歌。宦者贾举遮公从官而入。闭门。崔杼之徒持兵从中起。公登台而请解。不许。请盟。不许。请自杀于庙。不许。皆曰。﹁君之臣杼疾病。不能听命。近于公宫。陪臣争趣有淫者。不知二命。﹂公逾墙。射中公股。公反坠。遂弑之。晏婴立崔杼门外。曰。﹁君为社稷死则死之。为社稷亡则亡之。若为己死己亡。非其私昵。谁敢任之。﹂门开而入。枕公尸而哭。三踊而出。人谓崔杼。﹁必杀之。﹂崔杼曰。﹁民之望也。舍之得民。﹂\\
\hspace*{2em}丁丑。崔杼立庄公异母弟杵臼。是为景公。景公母。鲁叔孙宣伯女也。景公立。以崔杼为右相。庆封为左相。二相恐乱起。乃与国人盟曰。﹁不与崔庆者死。﹂晏子仰天曰。﹁婴所不获唯忠于君利社稷者是从。﹂不肯盟。庆封欲杀晏子。崔杼曰。﹁忠臣也。舍之。﹂齐太史书曰。﹁崔杼弑庄公。﹂崔杼杀之。其弟复书。崔杼复杀之。少弟复书。崔杼乃舍之。\\
\hspace*{2em}景公元年。初。崔杼生子成及强。其母死。取东郭女。生明。东郭女使其前夫子无咎与其弟偃相崔氏。成有罪。二相急治之。立明为太子。成请老于崔。崔杼许之。二相弗听。曰。﹁崔。宗邑。不可。﹂成。强怒。告庆封。庆封与崔杼有郄。欲其败也。成。强杀无咎。偃于崔杼家。家皆奔亡。崔杼怒。无人。使一宦者御。见庆封。庆封曰。﹁请为子诛之。﹂使崔杼仇卢蒲嫳攻崔氏。杀成。强。尽灭崔氏。崔杼妇自杀。崔杼毋归。亦自杀。庆封为相国。专权。\\
\hspace*{2em}三年十月。庆封出猎。初。庆封已杀崔杼。益骄。嗜酒好猎。不听政令。庆舍用政。已有内郄。田文子谓桓子曰。﹁乱将作。﹂田。鲍。高。栾氏相与谋庆氏。庆舍发甲围庆封宫。四家徒共击破之。庆封还。不得入。奔鲁。齐人让鲁。封奔吴。吴与之朱方。聚其族而居之。富于在齐。其秋。齐人徙葬庄公。僇崔杼尸于市以说众。\\
\hspace*{2em}九年。景公使晏婴之晋。与叔向私语曰。﹁齐政卒归田氏。田氏虽无大德。以公权私。有德于民。民爱之。﹂十二年。景公如晋。见平公。欲与伐燕。十八年。公复如晋。见昭公。二十六年。猎鲁郊。因入鲁。与晏婴俱问鲁礼。三十一年。鲁昭公辟季氏难。奔齐。齐欲以千社封之。子家止昭公。昭公乃请齐伐鲁。取郓以居昭公。\\
\hspace*{2em}三十二年。彗星见。景公坐柏寝。叹曰。﹁堂堂。谁有此乎。﹂群臣皆泣。晏子笑。公怒。晏子曰。﹁臣笑群臣谀甚。﹂景公曰。﹁彗星出东北。当齐分野。寡人以为忧。﹂晏子曰。﹁君高台深池。赋敛如弗得。刑罚恐弗胜。茀星将出。彗星何惧乎。﹂公曰。﹁可禳否。﹂晏子曰。﹁使神可祝而来。亦可禳而去也。百庶苦怨以万数。而君令一人禳之。安能胜众口乎。﹂是时景公好治宫室。聚狗马。奢侈。厚赋重刑。故晏子以此谏之。\\
\hspace*{2em}四十二年。吴王阖闾伐楚。入郢。\\
\hspace*{2em}四十七年。鲁阳虎攻其君。不胜。奔齐。请齐伐鲁。鲍子谏景公。乃囚阳虎。阳虎得亡。奔晋。\\
\hspace*{2em}四十八年。与鲁定公好会夹谷。犁锄曰。﹁孔丘知礼而怯。请令莱人为乐。因执鲁君。可得志。﹂景公害孔丘相鲁。惧其霸。故从犁锄之计。方会。进莱乐。孔子历阶上。使有司执莱人斩之。以礼让景公。景公惭。乃归鲁侵地以谢。而罢去。是岁。晏婴卒。\\
\hspace*{2em}五十五年。范。中行反其君于晋。晋攻之急。来请粟。田乞欲为乱。树党于逆臣。说景公曰。﹁范。中行数有德于齐。不可不救。﹂乃使乞救而输之粟。\\
\hspace*{2em}五十八年夏。景公夫人燕姬适子死。景公宠妾芮姬生子荼。荼少。其母贱。无行。诸大夫恐其为嗣。乃言愿择诸子长贤者为太子。景公老。恶言嗣事。又爱荼母。欲立之。惮发之口。乃谓诸大夫曰。﹁为乐耳。国何患无君乎。﹂秋。景公病。命国惠子。高昭子立少子荼为太子。逐群公子。迁之莱。景公卒。太子荼立。是为晏孺子。冬。未葬。而群公子畏诛。皆出亡。荼诸异母兄公子寿。驹。黔奔卫。公子驵。阳生奔鲁。莱人歌之曰。﹁景公死乎弗与埋。三军事乎弗与谋。师乎师乎。胡党之乎。﹂\\
\hspace*{2em}晏孺子元年春。田乞伪事高。国者。每朝。乞骖乘。言曰。﹁子得君。大夫皆自危。欲谋作乱。﹂又谓诸大夫曰。﹁高昭子可畏。及未发。先之。﹂大夫从之。六月。田乞。鲍牧乃与大夫以兵入公宫。攻高昭子。昭子闻之。与国惠子救公。公师败。田乞之徒追之。国惠子奔莒。遂反杀高昭子。晏圉奔鲁。八月。齐秉意兹。田乞败二相。乃使人之鲁召公子阳生。阳生至齐。私匿田乞家。十月戊子。田乞请诸大夫曰。﹁常之母有鱼菽之祭。幸来会饮。﹂会饮。田乞盛阳生橐中。置坐中央。发橐出阳生。曰。﹁此乃齐君矣。﹂大夫皆伏谒。将与大夫盟而立之。鲍牧醉。乞诬大夫曰。﹁吾与鲍牧谋共立阳生。﹂鲍牧怒曰。﹁子忘景公之命乎。﹂诸大夫相视欲悔。阳生前。顿首曰。﹁可则立之。否则已。﹂鲍牧恐祸起。乃复曰。﹁皆景公子也。何为不可。﹂乃与盟。立阳生。是为悼公。悼公入宫。使人迁晏孺子于骀。杀之幕下。而逐孺子母芮子。芮子故贱而孺子少。故无权。国人轻之。\\
\hspace*{2em}悼公元年。齐伐鲁。取讙。阐。初。阳生亡在鲁。季康子以其妹妻之。及归即位。使迎之。季姬与季鲂侯通。言其情。鲁弗敢与。故齐伐鲁。竟迎季姬。季姬嬖。齐复归鲁侵地。\\
\hspace*{2em}鲍子与悼公有郄。不善。四年。吴。鲁伐齐南方。鲍子弑悼公。赴于吴。吴王夫差哭于军门外三日。将从海入讨齐。齐人败之。吴师乃去。晋赵鞅伐齐。至赖而去。齐人共立悼公子壬。是为简公。\\
\hspace*{2em}简公四年春。初。简公与父阳生俱在鲁也。监止有宠焉。及即位。使为政。田成子惮之。骤顾于朝。御鞅言简公曰。﹁田。监不可并也。君其择焉。﹂弗听。子我夕。田逆杀人。逢之。遂捕以入。田氏方睦。使囚病而遗守囚者酒。醉而杀守者。得亡。子我盟诸田于陈宗。初。田豹欲为子我臣。使公孙言豹。豹有丧而止。后卒以为臣。幸于子我。子我谓曰。﹁吾尽逐田氏而立女。可乎。﹂对曰。﹁我远田氏矣。且其违者不过数人。何尽逐焉。﹂遂告田氏。子行曰。﹁彼得君。弗先。必祸子。﹂子行舍于公宫。\\
\hspace*{2em}夏五月壬申。成子兄弟四乘如公。子我在幄。出迎之。遂入。闭门。宦者御之。子行杀宦者。公与妇人饮酒于檀台。成子迁诸寝。公执戈将击之。太史子余曰。﹁非不利也。将除害也。﹂成子出舍于库。闻公犹怒。将出。曰。﹁何所无君。﹂子行拔剑曰。﹁需。事之贼也。谁非田宗。所不杀子者有如田宗。﹂乃止。子我归。属徒攻闱与大门。皆弗胜。乃出。田氏追之。丰丘人执子我以告。杀之郭关。成子将杀大陆子方。田逆请而免之。以公命取车于道。出雍门。田豹与之车。弗受。曰。﹁逆为余请。豹与余车。余有私焉。事子我而有私于其雠。何以见鲁。卫之士。﹂\\
\hspace*{2em}庚辰。田常执简公于俆州。公曰。﹁余蚤从御鞅言。不及此。﹂甲午。田常弑简公于俆州。田常乃立简公弟骜。是为平公。平公即位。田常相之。专齐之政。割齐安平以东为田氏封邑。\\
\hspace*{2em}平公八年。越灭吴。二十五年卒。子宣公积立。\\
\hspace*{2em}宣公五十一年卒。子康公贷立。田会反廪丘。\\
\hspace*{2em}康公二年。韩。魏。赵始列为诸侯。十九年。田常曾孙田和始为诸侯。迁康公海滨。\\
\hspace*{2em}二十六年。康公卒。卜氏遂绝其祀。田氏卒有齐国。为齐威王。强于天下。\\
\hspace*{2em}太史公曰。吾适齐。自泰山属之琅邪。北被于海。膏壤二千里。其民阔达多匿知。其天性也。以太公之圣。建国本。桓公之盛。修善政。以为诸侯会盟。称伯。不亦宜乎。洋洋哉。固大国之风也。\\

% ===== 索隐述赞（司马贞《史记索隐》）· 统一样式 =====
% 题头「索隐述赞曰」五字并入反白框；正文每四句一列，定宽盒 \hfil 居列中；
% \fontsize{24}{41} 列距对齐全书界行（121+41k）；\clearpage 使述赞独占一页。
\clearpage
\shuzanhead{索隐述赞曰}\linebreak
{\shuzanbody\fontsize{24}{41}\selectfont
\hbox to \hsize{\hfil 太公佐周。实秉阴谋。既表东海。乃居营丘。\hfil}\linebreak
\hbox to \hsize{\hfil 小白致霸。九合诸侯。及溺内宠。舋钟虫流。\hfil}\linebreak
\hbox to \hsize{\hfil 庄公失德。崔杼作仇。陈氏专政。原货轻收。\hfil}\linebreak
\hbox to \hsize{\hfil 悼简遘祸。田阚非俦。沨沨余烈。一变何由。\hfil}%
\par}
\end{document}
